\chapter{Curvature, Ricci, and the First Appearance of Matter}

Tonight we do not start over; we continue. The lecture resumes the curvature story from its operational beginning: carry a vector around a small closed loop and ask whether it comes back to itself. From there the discussion tightens in the two ways that matter for general relativity. First, the loop must be infinitesimal, because curvature is a local notion. Second, the geometry must be intrinsic, the sort of geometry that can be discovered by creatures or instruments living inside the space, not by peeking into an ambient embedding space. Only after those two points are in hand do we let the Riemann tensor appear. The lecture then makes a deliberate pivot: geometry is only half of gravity. The other half is matter, first rehearsed through charge and current, and then assembled into the stress-energy tensor. The lecture closes by pointing forward. Curvature belongs on one side of Einstein's equation, energy-momentum on the other, and geodesic motion tells us how matter responds.

\section{Recap: Curvature from Parallel Transport}

Let us begin where Susskind begins: with a two-dimensional surface. A straight line, incidentally, is just a special case of a curve, namely a curve that is uncurved. On a curved surface we still know what it means to move as straight as possible along the surface, and we still know what it means to parallel transport a vector. Curvature does not abolish these notions; it makes them nontrivial.

Take a vector tangent to the surface at some point of a closed path, and parallel transport it all the way around the path. If the region enclosed by the path is flat, the vector comes back to itself. If there is curvature in the enclosed region, the vector returns rotated by a small angle. That angle is the operational signal of curvature.

There are really two senses of angle to keep track of. One is the sense in which we circumnavigate the loop: clockwise or counterclockwise. The other is the sense in which the vector itself rotates. The sign of the curvature is determined by the relation between these two senses. If, for example, a counterclockwise excursion produces a counterclockwise rotation of the vector, we call the curvature positive. If a counterclockwise excursion produces a clockwise rotation, we call the curvature negative.

For very small loops the rotation angle is very small, and the curvature is defined by dividing that small deflection by the enclosed area and then taking the infinitesimal limit:
\begin{equation}
K=\lim_{\Delta A\to 0}\frac{\Delta\theta}{\Delta A}.
\end{equation}
This \(K\) is the curvature scalar of a two-dimensional surface.

Positive curvature is sphere-like; negative curvature is saddle-like. The lecture also gives the sphere as a benchmark example. If the surface is a sphere of radius \(a\), then
\begin{equation}
K_{\text{sphere}}=\frac{1}{a^2}.
\end{equation}
The larger the sphere, the smaller the curvature; the smaller the sphere, the larger the curvature. The important point is that we do not need to know where the center of the sphere is. The curvature is detected intrinsically, by local transport around a small loop.

\subsection{Question \& Answer}

Do these loops have to be small? Yes; in fact they must be taken infinitesimal. The lecture is quite explicit about this. Curvature is a local notion. We are not interested in heroic circumnavigations around the whole Earth. If our instruments were accurate enough, a sufficiently small loop would already reveal the curvature. That is why the definition is a limit as the enclosed area goes to zero. Large loops mix local curvature with global features of the surface; the differential-geometric quantity we want is obtained only in the small-loop limit.

\section{Intrinsic Versus Extrinsic Geometry}

The next interruption goes deeper. If we draw a curved two-dimensional surface, we almost automatically draw it embedded in a higher-dimensional space. But that is precisely the picture the lecture asks us to forget. We do not want to study the way the surface bends in its surroundings. We want to study what a creature living on the surface could discover without ever probing off the surface.

Susskind's bug is the right guide. The bug can see only within the surface, can lay rulers and protractors only within the surface, and can perform geometrical experiments only with tools attached to the surface. Whatever the bug can discover is intrinsic. By contrast, properties that depend on how the surface sits inside some larger space are extrinsic.

\subsection{Question \& Answer}

Are we measuring intrinsic curvature, or the curvature that comes from embedding the surface in a larger space? We are measuring the intrinsic curvature. That is the point of the bug thought experiment. The whole construction with parallel transport is meant to be something an inhabitant of the surface could test without ever looking in the perpendicular direction. In the language of the lecture, we want the geometry that a very small bug would discover by doing geometry on the surface itself.

This is why geodesics, parallel transport, and covariant differentiation belong together. They are intrinsic notions. A little wheeled device with its steering fixed straight ahead will move as straight as possible on the surface. On a sphere it traces a great circle. The lecture uses rollers first and then replaces them, more seriously, by gyroscopes. A gyroscope points only inside the space in which it lives. It does not point into some hypothetical external dimension. It therefore gives a physical model of intrinsic parallel transport.

Parallel transport in Riemannian geometry has another important property emphasized in the lecture: it preserves perpendicularity. If two vectors are perpendicular at the start, they remain perpendicular when transported. In two dimensions that almost fixes the notion by itself. In higher dimensions it is not quite enough, because a vector can remain perpendicular to a geodesic and still spiral around it. The lecture handles this by appealing to a family of neighboring geodesics. Two nearby geodesics define a plane, and one can require a transported vector to keep a fixed relation to that plane as well.

Operationally, the intuition is simple. On the Earth one keeps the steering wheel pointed dead straight ahead. One does not correct for mountains, passes, or embedding. One simply goes straight ahead locally. That is the geodesic idea in physical language. Technically, one chooses coordinates that are as nearly Cartesian as possible at the chosen point, so that the metric takes its flat form there and the Christoffel symbols vanish there. In that locally inertial frame a vector has no spurious coordinate-induced variation at the point itself. That is the mathematical expression of ``straight ahead.''

\section{From Small Loops to the Riemann Tensor}

Now the situation is clearly more complicated in higher dimensions. In three dimensions one may still speak loosely about an axis perpendicular to the little loop, but that shortcut is special to three dimensions. In general dimensions one must think in terms of planes. A plane is determined by two directions, and the curvature experiment involves two different planes.

The first is the plane of excursion: the infinitesimal plane in which we carry the vector around the loop. The second is the plane of rotation: the plane in which the vector is rotated when it returns. This is the real origin of the four indices of the curvature tensor.

So let us choose two small displacements, \(\delta x^\mu\) and \(dx^\nu\), spanning a tiny rectangle. The lecture itself corrects an early notational slip here: one must think of two distinct infinitesimal displacements, not two copies of the same \(dx\). The area enclosed is proportional to \(\delta x^\mu dx^\nu\). If we parallel transport a vector \(V^\rho\) around that loop, the change must vanish if there is no enclosed area, and therefore the leading nontrivial effect must be proportional to the area. It must also be linear in the starting vector. That is the geometric content behind the definition
\begin{equation}
\delta V^\rho = R^\rho{}_{\sigma\mu\nu}V^\sigma\,\delta x^\mu dx^\nu.
\end{equation}
Here we have written the canonical index arrangement rather than reproducing the lecture's intermediate false starts. The spoken discussion clearly identifies the roles of the indices, even though the notation is corrected on the fly.

\paragraph{A worked local update rule.}

The starting point is the transport equation along a curve:
\begin{equation}
\frac{D V^\rho}{ds}
=
\frac{dV^\rho}{ds}
+
\Gamma^\rho{}_{\sigma\mu}V^\sigma\frac{dx^\mu}{ds}
=
0.
\end{equation}
Multiplying by \(ds\) gives the infinitesimal change of a parallel-transported vector:
\begin{equation}
dV^\rho = -\Gamma^\rho{}_{\sigma\mu}V^\sigma\,dx^\mu.
\end{equation}
This is the first-order update rule. It says how the vector changes when we move a little bit in one direction. To go around the small rectangle, we apply this rule successively along two neighboring sides and compare the two possible orders. Because the Christoffel symbols vary from point to point, the two orders do not agree. Their difference is exactly the curvature effect. In canonical form the result is
\begin{align}
\delta V^\rho
&=
\left(
\partial_\mu \Gamma^\rho{}_{\nu\sigma}
-
\partial_\nu \Gamma^\rho{}_{\mu\sigma}
+
\Gamma^\rho{}_{\mu\lambda}\Gamma^\lambda{}_{\nu\sigma}
-
\Gamma^\rho{}_{\nu\lambda}\Gamma^\lambda{}_{\mu\sigma}
\right)
V^\sigma\,\delta x^\mu dx^\nu.
\end{align}
Therefore
\begin{equation}
R^\rho{}_{\sigma\mu\nu}
=
\partial_\mu \Gamma^\rho{}_{\nu\sigma}
-
\partial_\nu \Gamma^\rho{}_{\mu\sigma}
+
\Gamma^\rho{}_{\mu\lambda}\Gamma^\lambda{}_{\nu\sigma}
-
\Gamma^\rho{}_{\nu\lambda}\Gamma^\lambda{}_{\mu\sigma}.
\end{equation}

This formula is not introduced as something to memorize first and understand later. The lecture motivates it geometrically. Derivatives of the Christoffel symbols must appear because the Christoffel symbols vary from place to place. Quadratic terms must appear because transport is done in sequence and the successive linear maps multiply. The tensor is the mathematical probe that detects genuine curvature, not merely curvilinear coordinates.

There is one more technical worry that the lecture raises. If we try to form the tiny loop by following coordinate lines, the loop need not close perfectly. That is true. But the failure to close is itself already a second-order effect in the small displacements. It changes only terms of the same infinitesimal order as the area, and it does not disturb the leading curvature formula. The curvature is still read off from the first nontrivial mismatch after circumnavigating the little loop.

\begin{figure}[ht]
\centering
\includegraphics[width=0.72\textwidth]{lecture_07_figure_02.png}
\caption{Riemann tensor formula on the board.}
\end{figure}

\subsection{Question \& Answer}

Why is this really a tensor? The objection is a good one, and the lecture stops for it. A derivative of a vector is not generally a tensor, so why should the coefficient in the small-loop formula be tensorial?

The answer is that the left-hand side of the loop formula is not the derivative of vectors at different points. It is the difference between two vectors at the same point: the original vector and the vector obtained after going around the loop and returning to the starting point. A difference of two vectors at one point is again a vector. The displacements \(\delta x^\mu\) and \(dx^\nu\) are vectors, and \(V^\sigma\) is arbitrary. If
\begin{equation}
\delta V^\rho = R^\rho{}_{\sigma\mu\nu}V^\sigma\,\delta x^\mu dx^\nu
\end{equation}
holds for every starting vector \(V^\sigma\), then \(R^\rho{}_{\sigma\mu\nu}\) must transform as a tensor. That is the tensoriality argument.

The lecture also asks whether the second half of the formula is some kind of commutator. It is. If for each \(\mu\) we define a matrix by
\begin{equation}
(\Gamma_\mu)^\rho{}_\sigma=\Gamma^\rho{}_{\mu\sigma},
\end{equation}
then the quadratic part of the Riemann tensor is exactly the matrix commutator
\begin{equation}
(\Gamma_\mu\Gamma_\nu-\Gamma_\nu\Gamma_\mu)^\rho{}_\sigma
=
\Gamma^\rho{}_{\mu\lambda}\Gamma^\lambda{}_{\nu\sigma}
-
\Gamma^\rho{}_{\nu\lambda}\Gamma^\lambda{}_{\mu\sigma}.
\end{equation}
More abstractly, the whole curvature operator is the commutator of covariant derivatives:
\begin{equation}
[\nabla_\mu,\nabla_\nu]V^\rho = R^\rho{}_{\sigma\mu\nu}V^\sigma.
\end{equation}
Covariant derivatives in different directions do not commute, and the failure of commutation is curvature.

\section{Symmetries and the Ricci Contraction}

The lecture next turns from definition to cleanup. The raw Riemann tensor has four indices, but many of its components are related. The symmetries are most transparent after lowering the upper index:
\begin{equation}
R_{\lambda\sigma\mu\nu}=g_{\lambda\rho}R^\rho{}_{\sigma\mu\nu}.
\end{equation}
We then regard the all-lowered tensor as an object built from two index-pairs, each pair corresponding to one of the two planes in the problem.

The cleaned symmetries are
\begin{align}
R_{\alpha\beta\gamma\delta}
&=
-R_{\beta\alpha\gamma\delta}
=
-R_{\alpha\beta\delta\gamma},
\\
R_{\alpha\beta\gamma\delta}
&=
R_{\gamma\delta\alpha\beta}.
\end{align}
So the tensor is antisymmetric within each pair and symmetric under interchange of the two pairs. This is exactly what the geometry suggests.

The next move is contraction. We now build a two-index tensor from the four-index tensor by contracting one index that refers to the circumnavigation plane with one index that refers to the rotation plane. In the notation visible on the board,
\begin{equation}
R_{\beta\delta}=g^{\alpha\gamma}R_{\alpha\beta\gamma\delta}.
\end{equation}
This is the Ricci tensor. By the symmetries of the Riemann tensor, it is symmetric:
\begin{equation}
R_{\beta\delta}=R_{\delta\beta}.
\end{equation}

The lecture is very frank here. The Ricci tensor is not easy to picture. The Riemann tensor comes with a clean operational picture: transport a vector around a small loop and see how it fails to come back. Ricci comes from a contraction of that object, and its geometric meaning is far less immediate. But it is precisely this reduced tensor that will play the central dynamical role later.

The lecture also remarks on a useful three-dimensional simplification. In three dimensions a plane corresponds to an axis perpendicular to it, so an object that depends on two planes can be repackaged as an object that depends on two directions. That is why in three dimensions Ricci and Riemann carry the same information. In four-dimensional space-time that simplification no longer happens, and the contraction to Ricci is a genuine reduction.

\begin{figure}[ht]
\centering
\includegraphics[width=0.68\textwidth]{lecture_07_figure_03.png}
\caption{Ricci tensor as a contraction of the Riemann tensor.}
\end{figure}

\section{Local Conservation Warm-Up: Charge Density, Current, and Continuity}

At this stage the lecture deliberately changes gears. Geometry is only half of gravity; matter is the other half. Wheeler liked to summarize the theory by saying that mass curves geometry, and geometry tells particles how to move. The present lecture is building the first half of that slogan into mathematical form.

Susskind begins with charge because the bookkeeping is simpler there. Take a small box of ordinary space, not space-time yet, and ask how much charge it contains. If \(dq\) is the charge in a small volume \(dV\), then the charge density is
\begin{equation}
\rho_q=\frac{dq}{dV}.
\end{equation}
For a finite region \(V\), the total charge is
\begin{equation}
q(V)=\int_V \rho_q\,dV.
\end{equation}

Now turn from a box to a window. A current is the amount of charge passing through an oriented area per unit time. If the window is perpendicular to the \(x\)-direction, the corresponding current component is \(J_x\); similarly for \(J_y\) and \(J_z\). The lecture gives a nice dimensional reason for grouping density and flux together. With \(c=1\), time and length have the same units. Charge density is charge per three space-like directions. A current component is charge per two space-like directions and one time-like direction. They therefore have the same overall dimensional character and fit naturally together into a single four-vector:
\begin{equation}
J^\mu=(\rho_q,J_x,J_y,J_z).
\end{equation}

The lecture phrases this in terms of different kinds of space-time windows. A density corresponds to a space-space-space window. A current corresponds to a space-space-time window. There are no time-time-space windows because there is only one time direction. This is precisely why the components gather into a four-vector.

One point must be separated carefully from this. The total charge itself is a scalar. All observers agree about how much charge there is. The density and the current are not scalars; together they form the four-current. That distinction is a warm-up for the more delicate case of energy and momentum.

\subsection{Question \& Answer}

Why is global conservation not enough? Because global conservation by itself would be operationally useless. If charge could disappear here and reappear somewhere far away, the total charge of the universe might remain constant, but an experiment performed in this room would have every right to conclude that charge had failed to be conserved locally.

The real statement is local charge conservation: the only way the charge in a room can change is by observable flow through the walls. The mathematical object that measures net outward flow is the divergence. Therefore the change in charge density must be tied to the divergence of the current:
\begin{equation}
\partial_t\rho_q+\nabla\cdot\mathbf J=0.
\end{equation}
Written out in components,
\begin{equation}
\partial_t\rho_q+\partial_x J_x+\partial_y J_y+\partial_z J_z=0.
\end{equation}
If we now recognize \(\rho_q\) as the time component of \(J^\mu\), the continuity equation becomes
\begin{equation}
\partial_\mu J^\mu=0.
\end{equation}
This is the compact local form of charge conservation. It is the exact pattern the lecture wants us to carry over to energy and momentum.

\section{Energy, Momentum, and the Stress-Energy Tensor}

Charge was the easy case. Energy and momentum behave similarly in one respect and differently in another. Like charge, they are conserved. Unlike total charge, they are not invariants. Different observers assign different total energies and different total momenta to the same system. The object that transforms covariantly is the four-vector obtained by packaging them together:
\begin{equation}
P^\mu=(P^0,P^1,P^2,P^3)=(E,p_x,p_y,p_z).
\end{equation}
The surviving board frame at this moment catches only the beginning of the notation, so the full vector written above is reconstructed from the transcript in standard relativistic notation.

\begin{figure}[ht]
\centering
\includegraphics[width=0.52\textwidth]{lecture_07_figure_04.png}
\caption{Beginning of the four-momentum notation on the board.}
\end{figure}

But for field theory and for gravity, total quantities are not enough. We need local densities and flows. That is what leads from the four-vector \(P^\mu\) to the two-index tensor \(T^{\mu\nu}\).

The lecture builds the tensor row by row. The density of energy is called \(T^{00}\). The flow of energy along the \(x\)-direction is \(T^{01}\). The flow of energy along the \(y\)-direction is \(T^{02}\), and so on. Next comes the density of \(x\)-momentum, \(T^{10}\), then the flow of \(x\)-momentum along \(x\), \(T^{11}\), along \(y\), \(T^{12}\), and along \(z\), \(T^{13}\). The pattern then repeats for \(y\)-momentum and \(z\)-momentum.

The rule is simple and worth stating plainly:
\begin{itemize}
\item the first index tells us which conserved quantity we are talking about;
\item the second index tells us whether we mean density or the direction of flow.
\end{itemize}

Thus the sixteen components assemble into
\begin{equation}
T^{\mu\nu}
=
\left(
\begin{array}{cccc}
T^{00} & T^{01} & T^{02} & T^{03} \\
T^{10} & T^{11} & T^{12} & T^{13} \\
T^{20} & T^{21} & T^{22} & T^{23} \\
T^{30} & T^{31} & T^{32} & T^{33}
\end{array}
\right).
\end{equation}
The matrix looks formal, but the lecture's rhythm is the opposite of formalism for its own sake: each entry is introduced as an intelligible density or flow before it is wrapped into tensor notation.

\begin{figure}[ht]
\centering
\includegraphics[width=0.58\textwidth]{lecture_07_figure_05.png}
\caption{First stress-energy tensor components on the board.}
\end{figure}

A few entries deserve to be named explicitly:
\begin{align}
T^{00} &:\ \text{energy density}, \\
T^{01} &:\ \text{flow of energy along the }x\text{-direction}, \\
T^{10} &:\ \text{density of }x\text{-momentum}, \\
T^{12} &:\ \text{flow of }x\text{-momentum along the }y\text{-direction}.
\end{align}
The lecture also notes, without proof, a general property of relativistic systems: the energy-momentum tensor is symmetric,
\begin{equation}
T^{\mu\nu}=T^{\nu\mu}.
\end{equation}
So momentum density and energy flow are the same object viewed with the indices exchanged.

Now come the continuity equations. Each row of \(T^{\mu\nu}\) corresponds to a conserved quantity, and each row therefore has its own local continuity equation. For the energy row we have
\begin{equation}
\partial_t T^{00}+\partial_x T^{01}+\partial_y T^{02}+\partial_z T^{03}=0.
\end{equation}
For the \(x\)-momentum row,
\begin{equation}
\partial_t T^{10}+\partial_x T^{11}+\partial_y T^{12}+\partial_z T^{13}=0.
\end{equation}
And similarly for the remaining two rows. All four equations can be condensed into the single tensor statement
\begin{equation}
\partial_\mu T^{\nu\mu}=0.
\end{equation}

The lecture is again careful about the status of this equation. As written, it is the conservation law in special relativity. In general relativity the ordinary derivative will be replaced by a covariant derivative. That is a small-looking change with large consequences, and the lecture postpones it to the next step.

\section{Geodesic Motion and the Newtonian Preview}

By now the pieces are on the table. Geometry represents the gravitational field. Energy and momentum represent the source of that field. The task Einstein faced was to connect them by tensor equations that make sense in any frame.

But that is only half the theory. The other half is motion. What does geometry do to particles? Susskind states Einstein's principle in the simplest possible form: particles move on geodesics of space-time. Their steering wheels are locked straight ahead, not merely in space but in space-time.

In flat space-time this reduces to ordinary inertial motion. A straight line in space-time means a straight path in space together with constant velocity. In curved space-time the statement becomes
\begin{equation}
\frac{D}{d\tau}\left(\frac{dx^\mu}{d\tau}\right)=0,
\end{equation}
where \(\tau\) is proper time and \(dx^\mu/d\tau\) is the tangent vector to the worldline. Expanding the covariant derivative gives the familiar geodesic equation:
\begin{equation}
\frac{d^2x^\mu}{d\tau^2}
+
\Gamma^\mu{}_{\nu\sigma}
\frac{dx^\nu}{d\tau}
\frac{dx^\sigma}{d\tau}
=0.
\end{equation}
This is the differential equation for the orbit of a particle in a gravitational field.

The lecture ends with a preview of the Newtonian limit. Suppose the field is weak and all particles move slowly. Then the spatial components of \(dx^\mu/d\tau\) are small, while the time component is approximately
\begin{equation}
\frac{dt}{d\tau}\approx 1.
\end{equation}
Also, in this limit \(\tau\) and \(t\) are nearly the same, so the left-hand side of the spatial geodesic equation is nearly the ordinary acceleration. The dominant Christoffel term is the one with two time indices, and the spatial equation reduces to
\begin{equation}
\frac{d^2x^i}{dt^2}\approx -\Gamma^i{}_{00}.
\end{equation}
Thus, up to the sign already fixed by the geodesic equation, \(\Gamma^i{}_{00}\) plays the role of the Newtonian gravitational field, the field of acceleration. This is the bridge the lecture wants us to see: Newtonian gravity appears as the slow-motion shadow of geodesic motion in curved space-time.

There is a corresponding bridge on the source side. Mass is the Newtonian source of gravity, but with \(c=1\) we have \(E=m\), and in relativity energy cannot be separated from momentum. So once energy is admitted as a source, momentum must come along with it. That is why a tensor equation is inevitable. A single scalar source will not survive a change of frame.

\section{Summary}

The lecture unfolds in a deliberate order. We first rebuild curvature from the old two-dimensional picture of parallel transport around a small loop, insisting that the loop be infinitesimal and the geometry intrinsic. We then enlarge that picture to higher dimensions, where two planes must be tracked and the curvature becomes the four-index Riemann tensor. After cleaning up its tensor character and its commutator structure, we lower indices, extract the basic symmetries, and contract to the Ricci tensor. Only then do we pivot from geometry to matter. Charge and current provide the warm-up for local conservation, and energy and momentum then assemble into the stress-energy tensor with its four continuity equations. The chapter closes where the lecture closes: curvature on one side, stress-energy on the other, and geodesic motion as the rule that tells particles how to move through the geometry they help create.